\section{Related work}
\label{sec:related}

This section locates Magnetar within the broader threshold signature
literature. We organise by primitive family: (i) the sparse
threshold-hash-based-signature literature, (ii) the per-validator
aggregate-signature pattern for PQ consensus, (iii) the
lattice-family threshold schemes (Pulsar, Corona, Threshold Raccoon,
Hermine, Threshold ML-DSA), (iv) the classical Schnorr-family
threshold (FROST / RFC 9591), (v) threshold ECDSA (CGGMP21), and
(vi) the threshold-cryptography standards landscape (NIST IR 8214C).

\subsection{Threshold hash-based signatures}
\label{sec:related:hash}

Hash-based threshold signatures are the smallest sub-field of
threshold cryptography literature, both because the FIPS 205
standard is recent (2024) and because the underlying primitive has
no linear share-aggregation identity (\S\ref{sec:fullmpc:no-komlo}).

\paragraph{Cosmin--Khovratovich (SAC 2021) \cite{cosmin2021}.}
``Threshold SPHINCS+'' shares the SPHINCS+ secret via Shamir over
the hash output domain and reconstructs at sign time --- structurally
identical to Magnetar's reveal-and-aggregate appendix construction.
Their analysis is in the random-oracle model under the
aggregator-honest assumption, the same TCB-required posture
documented in Appendix~\ref{app:reveal}.

\paragraph{McGrew et al. (IACR ePrint 2019/793) \cite{mcgrew2019}.}
The selected-element threshold-HBS pattern --- parties hold shared
WOTS+ chain heads and FORS leaves, release shares only of the
selected elements per signature. This is the construction Magnetar
implements at \texttt{ref/go/pkg/thbs/} (Appendix~\ref{app:thbs}).
The v1 setup is dealer-backed; the open research path to
public-DKG is the subject of \S\ref{sec:fullmpc}.

\paragraph{Komlo--Stebila--Sullivan (MPTC 2025 draft).} A NIST MPTC
draft on distributed FORS evaluation for threshold SLH-DSA;
preliminary results align with the layer-(C) sign-time MPC
direction. Magnetar tracks this work as a candidate future
reference for a public-DKG THBS.

\paragraph{Buterin--Drake on hash-based BLS replacements.} Various
informal proposals (Ethereum research forums) have considered
SLH-DSA as a replacement for BLS in proof-of-stake consensus. The
arguments are about \emph{aggregation} more than \emph{threshold}:
non-interactive aggregation of $N$ SLH-DSA signatures is itself
non-trivial because the per-signature WOTS+ chains do not
aggregate, but \emph{collection} of $N$ signatures into a single
on-chain record is straightforward. Magnetar's per-validator
standalone primitive (\S\ref{sec:construction}) is exactly the
collection pattern, with the Z-Chain Groth16 rollup providing
constant-size compression on top.

\subsection{Per-validator aggregate signatures for PQ consensus}
\label{sec:related:per-validator-aggregate}

The per-validator-aggregate pattern Magnetar instantiates for
SLH-DSA is the same pattern already in production for ML-DSA in
the Lux Quasar stack.

\paragraph{Lux \texttt{QuasarCert.MLDSAProof}.} The Lux consensus
\texttt{QuasarCert} carries a parallel \texttt{MLDSAProof} field
collecting $K$ standalone ML-DSA signatures from quorum
validators. The wire shape is byte-equivalent to Magnetar's
\texttt{ValidatorAggregateCert}; the cryptographic primitive is
FIPS 204 ML-DSA \cite{fips204} rather than FIPS 205 SLH-DSA.
Magnetar's standalone primitive is the explicit hash-based
analog of this pattern; the two share the
$(\mathit{Signers}, \mathit{PubKeys}, \mathit{Sigs})$
parallel-slices wire shape and the Z-Chain Groth16 rollup
witness layout.

\paragraph{Ethereum SSZ aggregate-signature records.} The
Ethereum consensus spec carries per-validator BLS signatures
collected into \texttt{AggregateAndProof} records before
in-block aggregation. The collection pattern is the same as
Magnetar's; the cryptographic primitive (BLS) is classical and
admits non-interactive aggregation that hash-based primitives
do not. Magnetar substitutes the Z-Chain Groth16 rollup for
BLS's algebraic aggregation; the rollup achieves the same
constant-size on-chain footprint with a different cryptographic
backing.

\subsection{Lattice-family threshold signatures}
\label{sec:related:lattice}

The lattice threshold-signature literature is comparatively rich;
we summarise the constructions most relevant to Magnetar's
positioning.

\paragraph{Pulsar (Lux LP-073) \cite{lp073pulsarpaper}.}
The companion paper. Pulsar is a two-round Module-LWE threshold
signature over
$R_q = \mathbb{Z}_q[X]/(X^{256}+1)$ with $q \approx 2^{48}$. The
construction shares the master secret $\mathbf{s} \in
R_q^{n_{\mathrm{vec}}}$ via Shamir over the ring, and runs a
two-round MAC-authenticated protocol with linear share
aggregation. Pulsar's public-BFT-safe threshold mechanics
contrast with Magnetar's per-validator standalone path: Pulsar
distributes a single shared secret across the committee with
no party holding the master key; Magnetar keeps a separate
keypair per validator. Pulsar and Magnetar share the LSS
lifecycle adapter contract and the cSHAKE256 transcript
discipline.

\paragraph{Corona / Corona (Lux LP-104)
\cite{lp104corona,boschini2024corona}.}
The other companion. Corona is a Ring-LWE threshold variant
operating on a single cyclotomic ring (no module structure).
Same public-BFT-safe threshold mechanics as Pulsar. Corona's
trade-off vs. Pulsar: simpler ring structure (R-LWE rather than
M-LWE), larger signatures, slightly slower signing.

\paragraph{Threshold Raccoon (Eurocrypt 2024) \cite{thresholdraccoon}.}
Standard-lattice threshold signatures from MLWE/MSIS at $\sim 13$~KiB
signature size and $\sim 40$~KiB per-user communication. Targets up
to 1024 signers. Includes a DKG, identifiable abort, and non-
interactive aggregation. Compared to Magnetar:
\begin{itemize}[nosep,leftmargin=2em]
  \item Threshold Raccoon is \emph{lattice-based} (threshold);
        Magnetar is \emph{hash-based} (per-validator standalone).
        They serve different cryptographic-diversity roles.
  \item Threshold Raccoon's signature size at SHAKE-192-equivalent
        security is $\sim 13$~KiB, comparable to Magnetar's
        $\sim 16$~KB.
\end{itemize}

\paragraph{Hermine (NIST MPTC 2025) \cite{hermine}.} Threshold
lattice signatures with DKG, identifiable aborts, and proactive
refresh. A direct peer of Threshold Raccoon at the NIST MPTC
standardisation track. Same family-positioning trade-off versus
Magnetar.

\paragraph{Threshold ML-DSA (USENIX 2025) \cite{tmldsa2025}.} The
Celi--del~Pino--Espitau--Niot--Prest construction targets
threshold ML-DSA with masked shares and amortised rejection
sampling. The rejection-sampling circularity
(App.~A \cite{quasarcertsoundness}) and AGM/trapdoor assumptions
in the proof keep this construction outside the Lux Quasar
production scope. Lux uses per-validator standalone ML-DSA
signatures (the \texttt{QuasarCert.MLDSAProof} pattern of
\S\ref{sec:related:per-validator-aggregate}), complementary to
Magnetar's hash-based standalone role.

\paragraph{Two-round threshold ML-DSA \cite{boneh2018,damgard2021,
dttf2024}.} Boneh, Eskandarian, Kim, Shih sketched a generic
threshold construction over MLWE; Damg{\aa}rd, Orlandi, Takahashi,
Tibouchi produced the first concrete two-round threshold ML-DSA
with rejection-sampling correctness; subsequent work
\cite{dttf2024} extended this with algebraic one-more LWE
assumptions. These constructions share the lattice family with
Pulsar and Corona; Magnetar is orthogonal.

\subsection{Schnorr-family threshold (FROST)}
\label{sec:related:schnorr}

\paragraph{FROST / RFC 9591 \cite{frost,rfc9591,komlogoldberg20}.}
Two-round threshold Schnorr with linear share aggregation, the
elliptic-curve canonical reference for two-round threshold
protocols. RFC 9591 standardises the wire format for FROST over
secp256k1, P-256, and ristretto255. FROST's linear share-
aggregation identity $z = \sum_i z_i$ is the template that
Pulsar and Corona port to the lattice setting; Magnetar cannot
use this template because SLH-DSA has no analogous identity
(\S\ref{sec:fullmpc:no-komlo}).

FROST is \emph{classical} (broken by Shor's algorithm); it is the
classical fast-path lane in Lux's Quasar production stack
\cite{lp020}, complementary to the three PQ lanes (Pulsar,
Corona, Magnetar). The Lens kernel (Lux LP-103 \cite{lp103})
extends FROST with the same LSS lifecycle as the PQ lanes.

\subsection{Threshold ECDSA (CGGMP21)}
\label{sec:related:cggmp}

\paragraph{CGGMP21 \cite{cggmp21}.} The production ECDSA threshold
standard. Requires Paillier encryption for the multiplicative-to-
additive conversion and ships an extensive ZK proof suite to
handle malicious aborts. Complexity is approximately an order of
magnitude higher than FROST or Pulsar at the same $(t, n)$.

CGGMP21 and Magnetar are not redundant: CGGMP21 signs ECDSA-format
bridge transactions for Ethereum-compatible counterparties;
Magnetar signs PQ-format consensus blocks for the Lux Quasar
pipeline. The two coexist by jurisdiction in the Lux production
stack.

\subsection{BLS aggregate}
\label{sec:related:bls}

\paragraph{BLS12-381 \cite{bls,boneh2001short,lp075}.} The classical
aggregate-signature standard for blockchain consensus. The
aggregate is 48 bytes regardless of $K$; verification is one
elliptic-curve pairing. BLS is the classical fast-path lane in
Lux's Quasar production stack. It is not post-quantum --- Shor's
algorithm inverts ECDLP --- so the Polaris cert profile pairs it
with the three PQ lanes (Pulsar, Corona, Magnetar) to provide
post-compromise PQ insurance.

The structural pattern Magnetar implements --- $N$ per-validator
signatures collected at the consensus layer, compressed by a
rollup or aggregation primitive --- is the same pattern BLS uses
in production. The cryptographic primitive differs (BLS pairing
vs FIPS 205 hash-based), and the compression mechanism differs
(BLS algebraic aggregation vs Z-Chain Groth16 rollup over the
parallel-slices cert).

\subsection{Standards landscape}
\label{sec:related:standards}

\paragraph{NIST IR 8214C \cite{nistir8214c}.} The NIST Multi-Party
Threshold Cryptography (MPTC) project surveys threshold ECDSA,
threshold Schnorr, threshold ML-DSA, and threshold ML-KEM.
Threshold SLH-DSA is explicitly noted in IR 8214C \S5 as ``a
candidate for future addition.'' Magnetar's contribution to the
MPTC scope:
\begin{itemize}[nosep,leftmargin=2em]
  \item A production-grade reference implementation of the
        per-validator standalone primitive shipping FIPS 205
        byte-identity (the v0.5.2 artefact at
        \texttt{luxfi/magnetar} addresses the MPTC project's
        interest in concrete code artefacts).
  \item An honest disclosure of the threshold-form TEE
        requirement (Appendix~\ref{app:thbs},
        Appendix~\ref{app:reveal}), allowing the MPTC scope to
        be precise about what ``threshold SLH-DSA'' means in
        current practice (TEE-attested aggregator/dealer)
        versus what the open research direction targets
        (public-DKG with full MPC).
\end{itemize}

\paragraph{FIPS 205 \cite{fips205}.} The standard Magnetar
implements. The standard does not specify a threshold mode;
threshold constructions live outside the standard. Magnetar's
byte-identity claim binds against the FIPS 205 wire format and
the unmodified FIPS 205 verifier; this is what ``threshold
SLH-DSA'' should mean in deployment (the threshold or
collection layer is invisible to verifiers).

\subsection{Comparison table}
\label{sec:related:table}

Table~\ref{tab:comparison} summarises the structural comparison.

\begin{table}[H]
\centering
\small
\caption{Threshold / per-validator signature schemes compared.
``Sig size'' is the on-wire signature; ``Comm'' is per-party
broadcast bytes per signing session; ``Family'' identifies the
underlying mathematical structure; ``PQ'' marks post-quantum
security.}
\label{tab:comparison}
\begin{tabular}{lrlcccc}
\toprule
Scheme & Sig size & Family & Posture & DKG & Reshare & PQ \\
\midrule
\textbf{Magnetar standalone} (this paper) & 16 KB & hash-based & per-validator & N/A & per-key rotate & yes \\
\textbf{Magnetar THBS} (App.~\ref{app:thbs}) & 16 KB & hash-based & threshold (TEE) & dealer (v1) & v0.6+ & yes \\
\textbf{Magnetar reveal-and-aggregate} (App.~\ref{app:reveal}) & 16 KB & hash-based & threshold (TEE) & yes & v0.6+ & yes \\
Threshold SPHINCS+ \cite{cosmin2021} & 17 KB & hash-based & threshold & no & no & yes \\
Pulsar \cite{lp073pulsarpaper} & 33 KB & M-LWE       & threshold & yes & yes & yes \\
Corona \cite{lp104corona} & 56 KB & R-LWE          & threshold & yes & yes & yes \\
Threshold Raccoon \cite{thresholdraccoon} & 13 KiB & M-LWE & threshold & yes & no & yes \\
Hermine \cite{hermine} & 13 KiB & M-LWE           & threshold & yes & yes & yes \\
2R-Th-Dilithium \cite{damgard2021} & 30 KB & M-LWE   & threshold & no & no & yes \\
Threshold ML-DSA \cite{tmldsa2025} & 5 KB & M-LWE & threshold & no & no & yes${}^\star$ \\
FROST / RFC 9591 \cite{rfc9591} & 64 B & Schnorr   & threshold & yes & no${}^\dagger$ & no \\
CGGMP21 ECDSA \cite{cggmp21} & 64 B & ECDSA        & threshold & yes & yes & no \\
BLS aggregate \cite{bls} & 48 B & pairing          & $K$-of-$K$ aggregate & no & no & no \\
\bottomrule
\multicolumn{7}{l}{\footnotesize ${}^\star$ Threshold ML-DSA's PQ guarantee is
conditional on the rejection-sampling resolution; see App.~A
\cite{quasarcertsoundness}.} \\
\multicolumn{7}{l}{\footnotesize ${}^\dagger$ FROST refresh is provided by
the Lens kernel (LP-103 \cite{lp103}) via the LSS adapter
contract.}
\end{tabular}
\end{table}

\subsection{Where Magnetar's contribution lives}
\label{sec:related:contribution}

The cryptographic primitives in Magnetar are not new. FIPS 205
\cite{fips205} is the NIST-standardised SLH-DSA. The per-validator
standalone collection pattern is in production for ML-DSA via the
Lux \texttt{QuasarCert.MLDSAProof} field. The Shamir-VSS
reveal-and-aggregate pattern is folklore. The McGrew threshold-HBS
selected-element pattern is from 2019. What Magnetar contributes
is the \emph{composition}:
\begin{enumerate}[leftmargin=2em,nosep]
  \item A byte-identity-pinned per-validator standalone primary
        primitive that produces FIPS 205-conformant signatures
        validated empirically by
        $\textsc{TestMagnetar\_Wire\_FIPS205Verifiable}$ at all
        three SHAKE parameter sets.
  \item A canonical MAGS / MAGG wire codec siblinged with
        Pulsar's PULS / PULG and Corona's CORS / CORG framings,
        giving the consensus-layer dispatcher a uniform shape
        across the three Polaris lanes.
  \item Honest demotion of the threshold forms to TEE-only
        constructions for the M-Chain bridge custody and A-Chain
        confidential compute paths, with the public-BFT entry
        point pointing exclusively at the standalone primitive.
  \item An explicit-semantic standalone API
        (\texttt{PerValidatorKeypair}, \texttt{ValidatorSign},
        \texttt{VerifyAggregateCert}) that decomplects the
        public-BFT signing primitive from the threshold
        primitives that share the same Go package.
  \item A Polaris cert-profile composition that pairs the
        hash-based standalone lane with two lattice threshold
        lanes (Pulsar M-LWE, Corona R-LWE) to provide
        cryptographic-family diversity against any single-family
        algorithmic break.
\end{enumerate}

The crypto exists. The composition that turns it into the
hash-based leg of a three-lane Polaris cert profile is Magnetar's
contribution.

\section*{Acknowledgments}

The Lux Core Team acknowledges:
\begin{itemize}[nosep,leftmargin=2em]
  \item The FIPS 205 / SPHINCS+ authors \cite{sphincs,fips205}
        for the underlying primitive.
  \item The Cloudflare \texttt{circl} team \cite{circl} for the
        production-grade Go SLH-DSA reference that Magnetar
        wraps; the byte-identity claim binds against the
        \texttt{circl}-emitted FIPS 205 bytes.
  \item Cozzo, Smart, Bonte, and Tan for the threshold-HBS
        infeasibility analyses that grounded the architectural
        decision to use per-validator standalone for public BFT.
  \item McGrew, Fluhrer, Gazdag, Kampanakis, Morton, and
        Westerbaan for the selected-element threshold-HBS
        pattern (THBS) Magnetar instantiates for the TEE-only
        custody path.
  \item Shamir \cite{shamir1979} for the secret-sharing
        primitive used in the appendix threshold forms.
  \item The Pulsar (Lux LP-073) and Corona (Lux LP-104) authors
        for the companion lattice-family threshold constructions
        whose architectural pattern Magnetar's appendix forms
        inherit.
\end{itemize}
